% answers_short.tex —— 2006 年·北京卷·理科 · 简短答案速查
% 由 tools/make_short_answers.py 自 answers.tex 的 \gksolhead/\gkanswer 提取，
% 并与 2006/2006_数学真题_北京理_答案要点.md 逐题比对后生成，勿手改（改源 MD 后重跑）。

\begin{gknotice}
\textbf{使用说明}：本节只给最终答案便于快速核对，完整推导见解析版。答案为编写者独立推导，\textbf{非官方标准答案}。
\end{gknotice}

\gkpart{选择题}
\begin{center}\begin{tabular}{ccccccccc}
\toprule
题号 & 1 & 2 & 3 & 4 & 5 & 6 & 7 & 8 \\
\midrule
答案 & D & C & B & A & C & A & D & C \\
\bottomrule
\end{tabular}\end{center}

\gkpart{填空题}
\begin{itemize}[leftmargin=2.6em]
  \item[\textbf{9.}] $-\dfrac12$。
  \item[\textbf{10.}] $-14$。
  \item[\textbf{11.}] $\dfrac12$。
  \item[\textbf{12.}] $60^\circ$。
  \item[\textbf{13.}] 最小值 $\sqrt2$，最大值 $\sqrt{10}$。
  \item[\textbf{14.}] 球面距离 $\dfrac{\pi R}{3}$，球心到平面 $ABC$ 的距离 $\dfrac{\sqrt3R}{2}$。
\end{itemize}

\gkpart{解答题}
\begin{itemize}[leftmargin=2.6em]
  \item[\textbf{15.}] (1) $\{x\in\mathbb R:x\ne\frac\pi2+k\pi,\ k\in\mathbb Z\}$；(2) $\dfrac{14}{5}$
  \item[\textbf{16.}] (1) $x_0=1$；(2) $a=2,\ b=-9,\ c=12$
  \item[\textbf{17.}] (1) 证明见解析；(2) 证明见解析；(3) $135^\circ$
  \item[\textbf{18.}] (1) 方案一通过概率为 $ab+ac+bc-2abc$；方案二通过概率为 $\dfrac{ab+ac+bc}{3}$；(2) 方案一通过概率不小于方案二
  \item[\textbf{19.}] (1) $\dfrac{x^2}{2}-\dfrac{y^2}{2}=1$ 的右支（$x\ge\sqrt2$）；(2) 最小值为 $2$
  \item[\textbf{20.}] (1) 例如前十项为 $1,3,2,1,1,0,1,1,0,1$；(2) $a_n$ 的极限不存在；$b_n$ 的极限存在且为 $6$；(3) 证明见解析
\end{itemize}
